\section{Acquisition and automatic classification: from parsed document to
owned collection}
\label{sec:acquisition}

Acquisition determines what later retrieval can support, and it is where a
document is assigned to the collection that owns it. HMRS submits PDFs,
Office files and scans to MinerU through an asynchronous parsing service
\citep{mineru}; parse tasks return identifiers that callers poll to completion,
and OCR, layout and formula-extraction failures remain visible rather than
being stored as empty documents.

The intended ingest workflow follows observable stages: submit, parse,
agent-mediated quality review, duplicate decision, optional split, storage,
classification, indexing and a retrieval probe. Depending on the entry point,
artefacts may include a parse task identifier, a quality report, a content
fingerprint, a split manifest, a storage record, a placement decision and a
probe result. The workflow treats a content-derived retrieval probe as evidence
of retrievability; an HTTP success from a save call alone does not establish
that property. The platform accepts scanned pages, Office files and plain
text, but quality review, classification and completion checks are not all
enforced by the parse API.

\subsection{Quality gates and ingest lifecycle}

Duplicate handling has two paths rather than a universal admission gate. The
upload service rejects an exact same-parent duplicate by content hash, while
the separate near-duplicate tool reports similar candidates for review at a
configured threshold; it does not delete them automatically. In the ingest
workflow, an agent reviews parsed text for OCR noise, binary residue, bodyless
headings, whitespace collapse and language mismatch. Save and indexing behaviour
depends on the entry point: the web parsed-file save route stores content
without starting indexing, while the MCP parsed-document save operation can
start vector/BM25 indexing asynchronously. The ingestion agent classifies the document into a library whose scope matches
the parsed content's subject, not its filename, and the classification is
recorded so that the owning collection can be compared with the catalogue
afterwards.

When requested and above the configured context ceiling (30,000 characters in
the arXiv campaign; 60,000 in later configurations), documents are split at
recognised structural markers. The splitter targets 0.85 utilisation as a
soft target and detects units such as headings, Markdown tables, lists, block
quotes and fenced code. It does not guarantee preservation of equation
boundaries. Descriptions and tags are derived from the body and checked by the
ingest workflow; mismatches call for metadata correction, not document
rewriting.
Vector/BM25 indexing and graph construction are distinct operations.
Post-index verification checks selected properties such as a non-empty vector
collection and expected chunk identifiers; reconciliation can detect stale
index state, but is not a transaction across storage and index services.

\subsection{Acquisition evidence and accounting limits}

The 100-paper campaign produced 322 stored parts in five libraries, so the
collection level was populated across five distinct owning collections rather
than a single corpus.
Its reconciliation is incomplete: the state file records 91 submitted names;
at the 2026-10-13 evaluation snapshot 79 papers were retrievable and 21 were
not, and
the surviving artefacts do not enumerate attempts and failures per paper, so
the missing 21 cannot be attributed to parsing alone. Formula-dense scans and
a recovered batch-process crash were observed; neither establishes the cause
for every missing paper.

A separate cross-lingual test processed 400 Wikipedia pages per language:
ingestion success 1.0 in both rounds, and automatic membership accuracy 0.875
for English against 1.0 for Chinese and Japanese, which measures how often
content-based classification placed a document in the correct collection
independently of its filename, with the per-language storage and
self-retrieval results in Section~\ref{sec:e1}. These measurements
characterise ingestion, classification and storage, not answer quality.
